\chapter{Xây dựng bộ dữ liệu}
\label{ch:dulieu}

Không có bộ dữ liệu công khai nào được dựng sẵn cho tác vụ ``nhìn một màn hình, viết một
câu hướng dẫn cho người đọc''. Chương này trình bày cách chuyển AndroidControl thành tập
giám sát cho tác vụ đó. Toàn bộ quá trình \textbf{tự động hoàn toàn}: không thuê người
dán nhãn, và không có nhãn nào do một mô hình ngôn ngữ khác sinh ra nằm trong vòng huấn
luyện.

\section{Nguồn dữ liệu và phép ghép đa nguồn}

Mỗi bước cần ba thành phần: câu hướng dẫn do người viết, ảnh màn hình, và cây trợ năng.
Không bản phân phối công khai nào của AndroidControl mang đủ cả ba, nên phải ghép
(Bảng~\ref{tab:nguon}).

\begin{table}[t]
\caption{Hai bản phân phối trên Hugging Face và phần thiếu của mỗi bản.}
\label{tab:nguon}
\centering\small
\renewcommand{\arraystretch}{1.2}
\footnotesize
\begin{tabular}{@{}p{5.4cm}p{4.2cm}p{3.4cm}@{}}
\toprule
Kho & Có & Thiếu \\
\midrule
\code{HarrytheOrange/}\newline\code{\ \ parsed\_AndroidControl}
 & câu hướng dẫn đủ $15.283$ chuỗi & không có ảnh \\[0.4em]
\code{ckg/}\newline\code{\ \ AndroidControlParsed}\newline\code{\ \ WithImages-20k}
 & ảnh màn hình & mất câu hướng dẫn \\
\bottomrule
\end{tabular}
\end{table}

Hai kho được ghép theo khoá \code{(episode\_id, step\_id)}. Cây trợ năng của $99.131$ màn
lấy từ cùng bản phát hành gốc.

Một tính chất của bản phân phối phải nêu vì nó ảnh hưởng tới cách đọc kết quả: chuỗi mục
tiêu bị cắt ngay từ nguồn, $99{,}0\%$ kết thúc giữa chừng một cụm từ, ví dụ
\code{...look for Calvin Kle}. Mọi nhánh huấn luyện đều đọc cùng chuỗi mục tiêu bị cắt
đó, nên nó không giải thích được chênh lệch \emph{giữa} các nhánh; nhưng nhánh tham chiếu
(câu người viết) thì không thấy mục tiêu chút nào, nên khoảng cách giữa mô hình và câu
người viết ở chừng mực đó là một ước lượng thừa.

\section{Kiểm phép ghép}

Ghép lệch một bước là loại lỗi nguy hiểm nhất trong toàn bộ đường ống: số dòng vẫn đủ,
huấn luyện vẫn hội tụ, mọi thứ phía sau vẫn chạy trơn; chỉ có điều mô hình học ảnh của
bước $A$ với câu của bước $B$, và không khâu nào phía sau báo lỗi.

Phép kiểm được thiết kế như sau: chạy OCR tại điểm chạm đáp án, rồi xem chuỗi chữ đọc
được có xuất hiện trong câu hướng dẫn hay không. Nếu phép ghép đúng, tỉ lệ khớp phải cao
hơn hẳn so với một \emph{đối chứng lệch chủ ý một bước}. Trên $n = 400$ bước lấy mẫu:

\begin{center}
\begin{tabular}{@{}lr@{}}
\toprule
Ghép đúng & $48\%$ \\
Ghép lệch một bước (đối chứng) & $20\%$ \\
\midrule
Tỉ số & $2{,}4\times$ \\
\bottomrule
\end{tabular}
\end{center}

Nếu phép ghép sai thì hai con số phải xấp xỉ nhau; chúng cách nhau $2{,}4$ lần và khoảng
tin cậy không giao nhau. Tỉ lệ $48\%$ chỉ đọc được khi đặt cạnh $20\%$ của
đối chứng; mọi phép ghép về sau đều được kiểm theo cùng cách này.

\section{Quy mô, cách chia tập và kiểm rò rỉ}

\begin{table}[t]
\caption{Quy mô hai tập sau khi dựng.}
\label{tab:quymo}
\centering\small
\renewcommand{\arraystretch}{1.2}
\begin{tabular}{@{}lrr@{}}
\toprule
 & Tập dạy & Tập kiểm \\
\midrule
Số bước & $64.567$ & $6.958$ \\
Bước chạm (quần thể được chấm) & $41.191$ ($63{,}8\%$) & $4.463$ ($64{,}1\%$) \\
Số tác vụ & $12.895$ & $1.432$ \\
Độ phủ OCR & $100\%$ & $100\%$ \\
Độ phủ nhãn mô tả trên bước chạm & --- & $99{,}7\%$ \\
\bottomrule
\end{tabular}
\end{table}

Tập được chia \textbf{theo tác vụ}, không chia ngẫu nhiên theo dòng. Các màn hình trong
cùng một tác vụ gần như trùng nhau, nên chia ngẫu nhiên là rò rỉ trá hình và sẽ cho điểm
ảo. Phép kiểm rò rỉ được chạy ở quy mô đủ, đối chiếu $12.895$ tác vụ tập dạy với $1.432$
tác vụ tập kiểm, và cho kết quả \textbf{không có tác vụ nào trùng}. Đây là phép kiểm
không thể sửa được sau khi đã huấn luyện, nên nó được đặt làm một trong các chốt chặn
bắt buộc trước khi tiêu tài nguyên tính toán.

Ba tính chất của cách chia này giới hạn phạm vi kết luận và được trình bày đầy đủ.

\paragraph{Tách theo tác vụ không phải tách theo cách diễn đạt.} Trên một lát $1.697$
bước của tập dạy, tức $2{,}6\%$ tập dạy, đã có $17{,}3\%$ câu chuẩn của tập kiểm xuất
hiện \emph{nguyên văn} và $66{,}5\%$ chia sẻ ít nhất một cụm bốn từ, đơn giản vì bộ dữ
liệu dùng đi dùng lại các cách nói như ``click on the search bar''. Phép kiểm tái tạo ở
Mục~\ref{sec:phanbien} loại được những câu mà mô hình chép lại từ chính bước của nó,
nhưng không loại được hiện tượng này.

\paragraph{Tập kiểm không phải ``ứng dụng chưa từng thấy''.} Trong số các bước quy được
về một ứng dụng có tên, $95{,}6\%$ thuộc về ứng dụng cũng xuất hiện ở tập dạy
($2.991/3.130$ bước); tính theo ứng dụng riêng biệt thì tỉ lệ là $91{,}8\%$ ($247/269$). Phép so chính giữa
các nhánh không bị ảnh hưởng, vì mọi nhánh dùng chung dữ liệu dạy, chung mô hình gốc và
chung tập kiểm, nên chênh lệch giữa hai nhánh vẫn quy được về can thiệp. Tuy nhiên, một phép
so với một mô hình bên ngoài chưa từng thấy các ứng dụng này sẽ mang lợi thế sân nhà, và
luận văn không báo cáo phép so nào như vậy như một kết quả độc lập.
Mục~\ref{sec:latcat} đo trực tiếp ảnh hưởng của lợi thế này.

\paragraph{Hơn một nửa số bước không quy được về ứng dụng.} Chỉ $45{,}0\%$ số bước tập
kiểm gán được tên ứng dụng ($269$ ứng dụng); $3.828$ bước còn lại mang nhãn
\emph{không biết}. Nhãn này khác hẳn nhãn \emph{chưa thấy} và không bao giờ được đọc
thành \emph{chưa thấy}. Chính tính chất này quyết định luật gom cụm khi tính khoảng tin
cậy ở Mục~\ref{sec:thongke}.

\section{Trích chữ trên màn hình}

Cây trợ năng lẽ ra là nguồn tên phần tử tự nhiên, nhưng đo trên $99.131$ màn thì chỉ
$12{,}6\%$ phần tử mang bất kỳ tên nào, khớp với các báo cáo cho thấy phần lớn ứng
dụng Android xuất xưởng với nhãn trợ năng không đầy đủ~\cite{chen2020}. Ngoài ra, bản
dump hiện có chứa \code{content\_description} hướng về phía lập trình viên thay vì chữ
thật sự hiển thị trên màn.

Vì vậy toàn bộ ảnh màn hình được chạy qua OCR (RapidOCR, ONNX, chỉ dùng CPU), lấy
kết quả vừa làm đặc trưng đầu vào vừa làm nguồn tên phần tử, thu được $64.567$ bản ghi
không thiếu bản nào.

Khâu này được kiểm chéo giữa hai máy khác nhau, cách nhau năm ngày: $24{,}2$ chuỗi nhận
dạng được mỗi màn trên tập kiểm so với $23{,}6$ trên tập dạy, trung vị bằng nhau. Vì
AndroidControl không có nhãn chuẩn cho chữ trên màn nên luận văn \emph{không} tuyên bố
độ chính xác của OCR. Hai kiểu hỏng đã quan sát được: biểu tượng bị đọc thành ký tự, và
khoảng trắng bị nuốt (\code{Turnon}). Kiểu hỏng thứ hai đẩy phần tử vào nhóm ``không có
tên'' thay vì gán cho nó một tên sai, tức là nó làm nhãn nghèo đi chứ không làm nhãn
lệch đi.

\section{Nhãn mô tả tự động}
\label{sec:nhanmota}

Đích huấn luyện của nhánh đề xuất chứa một dòng khai báo có cấu trúc về phần tử đích.
Dòng này được dựng hoàn toàn bằng luật, không có người tham gia.

Quy trình: từ toạ độ chạm đáp án, chọn hộp trợ năng \emph{nhỏ nhất} chứa điểm đó; vai trò
lấy từ cây; tên lấy từ cây nếu vượt qua cổng lọc hợp lệ, ngược lại lấy từ OCR bên trong
hộp; ô toạ độ giữ chính điểm chạm đáp án quy về thang $[0, 1000]$; ô thứ tư giữ một dấu
hiệu phân biệt tính bằng luật.

\begin{table}[t]
\caption{Bốn ô của dòng khai báo và lý do tồn tại của từng ô.}
\label{tab:bono}
\centering\small
\renewcommand{\arraystretch}{1.25}
\footnotesize
\begin{tabular}{@{}p{2.4cm}p{4.3cm}p{5.9cm}@{}}
\toprule
Ô & Ví dụ & Vì sao cần \\
\midrule
Vai trò & \code{tappable text} & phân biệt nút, ô nhập, mục danh sách \\
Tên & \code{Search here} & chữ hiển thị; nguồn là OCR hoặc cây trợ năng \\
Toạ độ & \code{<point>440,91}\newline\code{</point>} & nhãn sạch $100\%$, chi phí dựng bằng không;
Bảng~\ref{tab:tienle} cho thấy bước trung gian không có toạ độ cho kết quả âm \\
Dấu hiệu phân biệt & \code{just above the text ``Coffee''} & thứ tách phần tử này khỏi
các phần tử giống nó; nhắm thẳng vào cơ chế hỏng ``câu mơ hồ'' \\
\bottomrule
\end{tabular}
\end{table}

\subsection{Hai cổng lọc được thêm vào sau khi đo}

Cả hai cổng dưới đây đều không có trong thiết kế ban đầu; chúng được thêm sau khi đo
thấy nhãn bị nhiễu.

\textbf{Cổng lọc nhãn trợ năng} loại bỏ $14{,}7\%$ nhãn tại các nút đích: định danh lấy
thẳng từ mã nguồn, chuỗi số dài, hậu tố dành cho trình đọc màn hình, và các cách diễn
đạt rỗng.

\textbf{Cổng hình học} chỉ cho phép lấy tên từ OCR khi hộp chiếm dưới $25\%$ diện tích
màn. Căn cứ: trong $25$ ca hộp lớn bằng gần cả màn hình, \emph{không ca nào} trong $25$
khớp với câu chuẩn.

\subsection{Ô dấu hiệu phân biệt phải làm lại}

Bản đầu của ô thứ tư gần như vô dụng: trên một lát thử, chỉ $7{,}3\%$ dấu hiệu thực sự
phân giải được điều gì, còn $85{,}8\%$ chỉ là đếm số lượng (``một trong bảy phần tử cùng
loại''), thứ không giúp người đọc hay mô hình phân biệt được gì. Sau khi sắp lại thứ tự
luật theo tiêu chí ``mệnh đề này có phân giải được sự mơ hồ không'' và bổ sung một mỏ neo
theo chữ (chuỗi OCR gần nhất nằm ngoài hộp kèm hướng tương đối), tỉ lệ này lên
$75{,}9\%$. Phần còn lại, $24{,}1\%$, vẫn không phân giải được gì và được nêu như một
giới hạn ở Mục~\ref{sec:hanche}.

\subsection{Chất lượng nhãn ở quy mô đủ}

Bảng~\ref{tab:chatluongnhan} đối chiếu nhãn dựng ở quy mô đủ ($41.191$ bước chạm) với
lát thử trước đó, nhỏ hơn $24$ lần.

\begin{table}[t]
\caption{Chất lượng nhãn mô tả ở quy mô đủ so với lát thử. Cả năm chỉ số lệch dưới một
điểm, cho thấy đường ống không suy giảm khi nhân quy mô.}
\label{tab:chatluongnhan}
\centering\small
\renewcommand{\arraystretch}{1.2}
\begin{tabular}{@{}lrr@{}}
\toprule
Chỉ số & Quy mô đủ & Lát thử \\
\midrule
Tên rõ ràng & $73{,}6\%$ & $73{,}8\%$ \\
Chỉ có ký hiệu, không có chữ & $4{,}4\%$ & --- \\
Không có tên & $22{,}0\%$ & $23{,}4\%$ \\
Vai trò rõ ràng & $75{,}8\%$ & --- \\
Trùng tên trên cùng màn & $7{,}6\%$ & $6{,}9\%$ \\
Có phần tử cùng vai trò để so & $91{,}6\%$ & --- \\
\bottomrule
\end{tabular}
\end{table}

Trong $32.061$ tên lấy được, $8.581$ đến từ cây trợ năng và $23.480$ đến từ OCR; tỉ lệ
này là hệ quả trực tiếp của việc chỉ $12{,}6\%$ phần tử có nhãn trợ năng.

Hai lỗi bắt được trong quá trình dựng đáng ghi lại vì cùng một loại: cây trợ năng lồng
nhau khiến chỉ số ``trùng tên'' bị thổi từ $30\%$ xuống giá trị thật $6{,}9\%$, do cùng
một nút bị đếm nhiều lần ở nhiều tầng của cây; và tên lớp Android
(\code{android.widget.TextView}) lọt vào ô tên phần tử. Cả hai đều thuộc loại lỗi không làm chương
trình dừng: kết quả vẫn ra và vẫn hợp lý, chỉ có giá trị là sai.

\subsection{Một quyết định đã tranh luận rồi bác bỏ}

Có giai đoạn tưởng phải đảo hẳn thứ tự ưu tiên sang lấy tên từ OCR trước, vì câu chuẩn
khớp OCR gấp $3{,}2$ lần khớp nhãn trợ năng. Truy lại cho thấy phạm vi ảnh hưởng thật chỉ
là $72/1074 = 6{,}7\%$ số ca; bỏ tiếp $32$ ca mà chuỗi số vốn đã có sẵn trong câu mục
tiêu thì còn $40$ ca lõi (OCR thắng $12$, trợ năng thắng $3$); và trong $12$ ca OCR
thắng thì $11$ ca là do nhãn trợ năng vốn đã là rác. Nghĩa là ưu thế của OCR là ưu thế
của \emph{việc lọc rác}, không phải của \emph{thứ tự ưu tiên}. Kết luận: giữ nguyên thứ
tự, thêm cổng lọc. Tỉ số $3{,}2$ lần vì vậy chỉ áp lên $6{,}7\%$ số ca, và trong phần đó
phần lớn quy về chất lượng nhãn chứ không quy về thứ tự ưu tiên.

\section{Bốn nhánh dữ liệu và các bất biến cấu trúc}
\label{sec:bonnhanh}

Bốn tập dữ liệu được dựng từ \emph{cùng một đầu vào} và chỉ khác nhau ở đích sinh. Lấy
đúng một bước thật (chuỗi $5590$, bước $1$) làm ví dụ:

\begin{description}[leftmargin=1.6em,style=nextline,font=\normalfont\bfseries]
\item[s1: câu trơn, mốc so nội bộ]
{\footnotesize\ttfamily Click on the search bar}

\item[s2: khai báo thật rồi tới câu]
{\footnotesize\ttfamily <desc>tappable text | Search here | <point>440,91</point>\\
\hspace*{1em}| just above the text ``Coffee''</desc>\\
Click on the search bar}

\item[s2r: khai báo GIẢ lấy từ màn khác, toạ độ ngẫu nhiên, khớp độ dài token]
{\footnotesize\ttfamily <desc>item | CATEGORIES | <point>713,402</point>\\
\hspace*{1em}| to the left of the text ``MEN''</desc>\\
Click on the search bar}

\item[s2\_nopoint: bỏ ô toạ độ, giữ ba ô còn lại]
{\footnotesize\ttfamily <desc>tappable text | Search here\\
\hspace*{1em}| just above the text ``Coffee''</desc>\\
Click on the search bar}
\end{description}

Nhánh \textbf{s2r} là nhánh đối chứng chưa có tiền lệ nào chạy. Nó chặn đúng cách giải
thích thay thế nguy hiểm nhất cho một kết quả dương của s2: ``chuỗi đích dài thêm thì mô
hình học kỹ hơn, chẳng liên quan gì tới nội dung mô tả''. Để nó chặn được, độ dài khai
báo giả phải sát độ dài thật. Cách ghép cuối cùng là lấy $12$ bản ghi gần nhất theo
\emph{độ dài token} rồi bốc ngẫu nhiên một bản theo hạt giống cố định, đạt $99{,}3\%$ số
cặp nằm trong hai token. Bản đầu ghép theo \emph{ký tự} cho trung vị lệch $0$ ký tự; đo lại theo token thì chỉ $54\%$ số cặp nằm trong hai token với
biên độ tới $\pm17$, mà hàm mất mát tính trên token chứ không tính trên ký tự.

\paragraph{Chín bất biến.} Sau mỗi lần dựng lại dữ liệu, chín bất biến cấu trúc được
kiểm và cả chín đều đạt ở quy mô đủ:

\begin{enumerate}[leftmargin=1.6em,itemsep=0.1em]
\item câu được chấm giống hệt nhau từng byte ở mọi nhánh, trên cả $64.567$ mẫu;
\item nhánh s1 không mang dòng khai báo nào;
\item nhánh s2\_nopoint không mang ô toạ độ nào;
\item ba nhánh có khai báo mang cùng $41.099$ nhãn;
\item s2 khác s2r đúng ở các nhãn đó và không khác ở chỗ nào khác;
\item khai báo giả không bao giờ tái sử dụng toạ độ đích thật;
\item các bước không phải bước chạm mang đích y hệt nhau ở cả bốn nhánh;
\item khoá ghép và số dòng khớp nhau giữa các nhánh;
\item độ dài token của khai báo thật và khai báo giả khớp theo ngưỡng đã nêu.
\end{enumerate}

Bất biến số $7$ là điều kiện tiên quyết cho phép kiểm không-gây-hại: vì $36{,}2\%$ dữ
liệu là bước không chạm (cuộn, gõ, mở ứng dụng, quay lại) và phần dữ liệu đó giống hệt
nhau ở bốn nhánh, nếu nhánh s2 làm hỏng nhóm bước này thì thấy ngay và không đổ được cho
sự khác biệt dữ liệu.
